\section{符号说明}

为便于阅读，将贯穿全文的核心符号列于表\ref{tab:notation}。各问题特有的局部符号在首次出现时另行定义。

\begingroup  
\rowcolors{2}{gray!10}{white}   

\begin{table}[!htbp]
\centering
\caption{核心符号及其含义}
\label{tab:notation}
\begin{tabularx}{\textwidth}{c X}
\toprule
\rowcolor{CumcmLightBlue}
符号 & 含义 \\
\midrule
$i$ & 任务索引，$i=1,2,\dots,50000$ \\
$r$ & 区域索引，$r\in\{\text{A},\text{B},\text{C},\text{D},\text{E},\text{F}\}$ \\
$t$ & 小时索引，$t=0,1,\dots,2406$ \\
$k_i$ & 任务类型，$k_i\in\{\text{RT},\text{Batch},\text{Training}\}$ \\
$a_i$ & 任务到达时刻（小时） \\
$p_i$ & 任务持续时长（小时） \\
$g_i$ & 任务GPU需求数（个） \\
$d_i$ & 任务最晚完成时刻（小时） \\
$L_i$ & 任务最大允许网络时延（ms） \\
$o_i$ & 任务来源区域 \\
$s_i$ & 任务开工时刻（小时） \\
$x_{i,r,s}$ & 二元决策变量：任务$i$在区域$r$、时刻$s$开工取1，否则取0 \\
$\Omega_i$ & 任务$i$的完整合法时空候选集 \\
$R_i$ & 任务$i$的空间可行域（满足网络时延约束的区域集合） \\
$T_i$ & 任务$i$的时间可行域（满足截止时间约束的开工时刻集合） \\
$\text{Slack}_i$ & 任务$i$的时间余量，$\text{Slack}_i = \min(d_i, 2406) - a_i - p_i$ \\
$P_{\text{AI}}[r,t]$ & 区域$r$在小时$t$的AI IT功率（MW） \\
$L[r,t]$ & 区域$r$在小时$t$的设施总负荷（MW），$L=\text{PUE}\times(\text{AI\_IT}+\text{NonAI\_IT})$ \\
$\text{PUE}[r]$ & 区域$r$的电源使用效率 \\
$\text{AvailableGPU}[r]$ & 区域$r$的可用GPU容量（个） \\
$\text{MaxITPower}[r]$ & 区域$r$的最大IT功率（MW） \\
$\text{MaxFacilityPower}[r]$ & 区域$r$的最大设施功率（MW） \\
$\text{GridPurchase}[r,t]$ & 区域$r$在小时$t$的电网购电功率（MW） \\
$\text{GridSell}[r,t]$ & 区域$r$在小时$t$的新能源售电功率（MW） \\
$\text{Curtailment}[r,t]$ & 区域$r$在小时$t$的新能源弃电功率（MW） \\
$\text{AvailableRenewable}[r,t]$ & 区域$r$在小时$t$的可用新能源功率（MW） \\
$\text{Price}[r,t]$ & 区域$r$在小时$t$的购电价格（CNY/MWh） \\
$\text{SellPrice}[r,t]$ & 区域$r$在小时$t$的售电价格（CNY/MWh） \\
$\text{CarbonIntensity}[r,t]$ & 区域$r$在小时$t$的电网碳强度（tCO$_2$/MWh） \\
$E[r,t]$ & 区域$r$在小时$t$末的储能荷电状态（SOC，MWh） \\
$\eta_c[r]$、$\eta_d[r]$ & 区域$r$的储能充电、放电效率 \\
$\text{StorageCapacity}[r]$ & 区域$r$的储能容量（MWh） \\
$\text{MaxChargePower}[r]$ & 区域$r$的最大充电功率（MW） \\
$\text{MaxDischargePower}[r]$ & 区域$r$的最大放电功率（MW） \\
$\text{MaxGridImport}[r]$ & 区域$r$的最大电网购电功率（MW） \\
$\text{MaxGridExport}[r]$ & 区域$r$的最大电网售电功率（MW） \\
$\text{SellLimit}[r]$ & 区域$r$的新能源售电限额（MW） \\
$N_t$ & 第$t$小时新到达任务总数 \\
$\lambda$ & 小时平均到达率（任务数/小时） \\
$\pi_{rk}$ & 任务源区域$r$与任务类型$k$的联合业务概率 \\
$\alpha_k$ & 任务类型$k$的单位GPU IT功率（MW/GPU） \\
$\omega(i,s,t)$ & 任务$i$在小时$t$与执行区间$[s,s+p_i)$的重叠时长（小时） \\
$\Delta L[r,t]$ & 区域$r$在小时$t$的设施负荷增量（MW） \\
$\Delta G[r,t]$ & 区域$r$在小时$t$的电网购电增量（MW） \\
$Q_r(L_r)$ & 区域$r$给定负荷曲线下的能源recourse最优值（CNY） \\
$\theta_r$ & Benders分解中区域$r$的能源成本代理变量（CNY） \\
\bottomrule
\end{tabularx}
\end{table}

\endgroup 